% from ega4-12-fr.tex
\oldpage[IV-3]{185}
sont affines, puis au cas où $S$ est \emph{noethérien}~; en effet, par (8.9.1)
et (8.5.2), on sait qu'il existe un préschéma noethérien
$S'=\operatorname{Spec}(A')$, où $A'$ est un sous-anneau de $A$, un morphisme
de type fini $f':X'\to S'$ et trois $\mathcal O_{X'}$-Modules cohérents
$\mathcal L'_i$ ($0\leq i\leq2$) tels que $X=X'\times_{S'}S$,
$f=f'_{(S)}$, $\mathcal L_i=\mathcal L'_i\otimes_{S'}S$, ainsi que deux
homomorphismes $u'$, $v'$ tels que $u=u'\otimes1$, $v=v'\otimes1$ et
$v'\circ u'=0$. On peut en outre supposer que $\mathcal L'_0$ est $f'$-plat
(11.2.7), et que $\mathcal L'_1$ est $f'$-plat lorsque $\mathcal L_1$ est
supposé $f$-plat. Par fidèle platitude, l'hypothèse
$(\mathcal H_n((\mathcal L_\bullet)_{f(x)}))_x=0$ équivaut à
$(\mathcal H_n((\mathcal L'_\bullet)_{f'(x')}))_{x'}=0$, si $x'$ est la
projection de $x$ dans $X'$~; si $U'$ est l'ensemble des $x'\in X'$ tels que
$(\mathcal H_n((\mathcal L'_\bullet)_{f'(x')}))_{x'}=0$, $U$ est donc l'image
réciproque de $U'$ par la projection $X\to X'$, ce qui ramène, pour la première
assertion, au cas où $S$ est noethérien. Pour la seconde assertion, on remarque
en outre que $A$ est limite inductive de ses sous-anneaux $A_\lambda$ qui sont
des $A'$-algèbres de type fini~; posons
$X_\lambda=X'\times_{S'}S_\lambda$, où
$S_\lambda=\operatorname{Spec}(A_\lambda)$,
$\mathcal L_i^{(\lambda)}=\mathcal L'_i\otimes_{S'}S_\lambda$, et soient
$u_\lambda=u'\otimes1:\mathcal L_2^{(\lambda)}\to\mathcal L_1^{(\lambda)}$,
$v_\lambda=v'\otimes1:\mathcal L_1^{(\lambda)}\to\mathcal L_0^{(\lambda)}$,
de sorte que l'on a
$\mathcal H_1(\mathcal L_\bullet^{(\lambda)})
=\operatorname{Ker}(v_\lambda)/\operatorname{Im}(u_\lambda)$. Or, comme le
foncteur $\varinjlim$ est exact dans la catégorie des groupes commutatifs, on a
$\operatorname{Ker}(v)=\varinjlim\operatorname{Ker}(v_\lambda)$,
$\operatorname{Im}(u)=\varinjlim\operatorname{Im}(u_\lambda)$, et
$\mathcal H_1(\mathcal L_\bullet)
=\varinjlim\mathcal H_1(\mathcal L_\bullet^{(\lambda)})$. Si l'on a supposé
que $\mathcal L_1$ est $f$-plat et (en se ramenant au cas où $X=U$) que l'on
ait prouvé $\mathcal H_1(\mathcal L_\bullet^{(\lambda)})=0$ pour tout
$\lambda$, on en déduira bien $\mathcal H_1(\mathcal L_\bullet)=0$.

\medskip
I) Supposons désormais $S$ noethérien. On sait (sans hypothèse de platitude sur
$\mathcal L_0$) que l'ensemble $U$ est \emph{constructible} dans $X$ (9.9.6).
Utilisant maintenant ($\mathbf 0_{\mathrm{III}}$, 9.2.5), il reste à montrer
que pour toute \emph{générisation} $x'$ de $x\in U$ dans $X$, on a aussi
$x'\in U$. La méthode exposée dans (12.0.2) s'applique sans changement (compte
tenu du fait que pour un préschéma $Z$ sur un corps $k$ et une extension $k'$
de $k$, la projection $Z\otimes_k k'\to Z$ est fidèlement plate). On peut donc
supposer que $S$ est le spectre d'un anneau de valuation discrète $A$, dont on
désigne par $t$ une uniformisante, $x$ étant au-dessus du point fermé de $S$ et
$x'$ au-dessus du point générique de $S$. L'hypothèse que $\mathcal L_0$ est
$f$-plat entraîne que $t$ est $\mathcal L_0$-régulier
($\mathbf 0_{\mathrm I}$, 6.3.4). On est alors ramené à prouver le lemme suivant
(où $B$, $M$, $N$, $P$ seront remplacés par $\mathcal O_x$,
$(\mathcal L_2)_x$, $(\mathcal L_1)_x$ et $(\mathcal L_0)_x$ respectivement)~:

\begin{lemma}[12.3.3.1]
\label{IV.12.3.3.1-fr}
Soient $B$ un anneau local noethérien, $M,N,P$ trois $B$-modules, $N$ étant
supposé de type fini, $M\xrightarrow{u}N\xrightarrow{v}P$ deux homomorphismes
tels que $v\circ u=0$. Soit $t$ un élément de l'idéal maximal de $B$ tel que
$t$ soit $P$-régulier et que la suite
\begin{equation}
  \tag{12.3.3.2}
  \label{IV.12.3.3.2-fr}
  M/tM\xrightarrow{u\otimes1_{B/tB}}N/tN
  \xrightarrow{v\otimes1_{B/tB}}P/tP
\end{equation}
soit exacte. Alors la suite $M\xrightarrow{u}N\xrightarrow{v}P$ est exacte.
\end{lemma}

Notons d'abord que si l'on remplace $M$ par son image $Q$ dans $N$ et $u$ par
l'injection $j:Q\to N$, l'image de $j\otimes1$ est la même que celle de
$u\otimes1$, donc on peut, sans changer ni l'hypothèse ni la conclusion,
supposer $u$ injectif. D'autre part, si $R$ est l'image de $v$ et
$\rho:N\to R$ la surjection canonique, $t$ est évidemment $R$-régulier, on a
$\rho\circ u=0$, et le noyau de $\rho\otimes1$ est contenu dans celui de
$v\otimes1$~; il lui est par suite égal si la suite (12.3.3.2) est exacte~;
comme d'autre part on a évidemment $\operatorname{Ker}(\rho)
=\operatorname{Ker}(v)$,
